\begin{longtable}{|>{\scriptsize\raggedright\arraybackslash}p{3cm}|>{\raggedright\arraybackslash}p{4cm}|>{\raggedright\arraybackslash}p{5.5cm}|>{\raggedright\arraybackslash}p{5.5cm}|}
\caption{Integration Test Cases}\label{tab:integration_test_cases}\\
\hline
\textbf{Test Case ID} & \textbf{Test Title / Scenario} & \textbf{Test Steps} & \textbf{Expected Results} \\
\hline
\endfirsthead
\multicolumn{4}{c}%
{{\bfseries \tablename\ \thetable{} -- continued from previous page}} \\
\hline
\textbf{Test Case ID} & \textbf{Test Title / Scenario} & \textbf{Test Steps} & \textbf{Expected Results} \\
\hline
\endhead
\hline \multicolumn{4}{|r|}{{Continued on next page}} \\ \hline
\endfoot
\hline
\endlastfoot

% Scenario 1: Full Book and User Lifecycle
\multicolumn{4}{|l|}{\textbf{Scenario: Full Lifecycle - From Entity Creation to Borrowal}} \\ \hline
TC\_INT\_001 & End-to-end workflow of creating all required entities for a new book, registering a new member, and the member borrowing that book. & 
    \begin{itemize}[noitemsep, topsep=0pt, partopsep=0pt, leftmargin=*]
        \item \textbf{As Librarian:} Login to the system.
        \item Navigate to Author Management and add a new author (e.g., "J.R.R. Tolkien").
        \item Navigate to Genre Management and add a new genre (e.g., "Fantasy").
        \item Navigate to Book Management and add a new book ("The Hobbit"), linking the newly created author and genre. Set `isAvailable` to true.
        \item Navigate to User Management and register a new user with `isAdmin=false` (e.g., "Member1").
        \item Logout.
        \item \textbf{As Member1:} Login to the system.
        \item Navigate to the book list and find "The Hobbit".
        \item Initiate and confirm the borrowing of "The Hobbit".
    \end{itemize} & 
    \begin{itemize}[noitemsep, topsep=0pt, partopsep=0pt, leftmargin=*]
        \item The author is created successfully and selectable in the book form.
        \item The genre is created successfully and selectable in the book form.
        \item The book is created with correct author and genre associations. The book list shows its status as available.
        \item The new member account is created and appears in the user list.
        \item Member1 can log in successfully and is directed to the member landing page.
        \item A new borrowal record is created for Member1 and "The Hobbit".
        \item The `isAvailable` status of "The Hobbit" in the Book entity is updated to `false`.
        \item The new borrowal record appears in Member1's borrowal history.
    \end{itemize} \\ \hline

% Scenario 2: Data Integrity on Deletion
\multicolumn{4}{|l|}{\textbf{Scenario: Data Integrity and Referential Constraints}} \\ \hline
TC\_INT\_002 & Verify data integrity when a linked entity (Author) is deleted after being associated with another entity (Book). & 
    \begin{itemize}[noitemsep, topsep=0pt, partopsep=0pt, leftmargin=*]
        \item \textbf{As Librarian:} Login.
        \item Create a new, unique Author (e.g., "IntegTest Author").
        \item Create a new Book (e.g., "IntegTest Book") and link it to "IntegTest Author".
        \item Verify the book details page shows "IntegTest Author" as the author.
        \item Navigate to Author Management and delete "IntegTest Author".
        \item Navigate back to Book Management and view the details of "IntegTest Book".
    \end{itemize} & 
    \begin{itemize}[noitemsep, topsep=0pt, partopsep=0pt, leftmargin=*]
        \item The Author is successfully deleted.
        \item The Book "IntegTest Book" still exists in the system.
        \item The `authorId` field for "IntegTest Book" is handled gracefully (e.g., set to null or an empty value), as the data model indicates it's optional. The book details page no longer displays the deleted author's name.
        \item The system does not crash and remains in a consistent state.
    \end{itemize} \\ \hline
TC\_INT\_003 & Verify that deleting a User with active Borrowals does not violate data consistency. &
    \begin{itemize}[noitemsep, topsep=0pt, partopsep=0pt, leftmargin=*]
        \item \textbf{As Librarian:} Login.
        \item Ensure a specific Member (e.g., "MemberToDelete") has at least one active borrowal record. If not, create one.
        \item Navigate to User Management.
        \item Attempt to delete the user "MemberToDelete".
    \end{itemize} & 
    \begin{itemize}[noitemsep, topsep=0pt, partopsep=0pt, leftmargin=*]
        \item \textbf{Expected (Option A - Deletion Blocked):} The system prevents the deletion and displays an error message indicating the user has active borrowals that must be resolved first.
        \item \textbf{Expected (Option B - Soft Delete/Orphaning):} The user is deleted. The `memberId` in the Borrowal record  remains but now points to a non-existent user. The borrowal record is still visible to the Librarian, but may indicate the user is deleted.
        \item (The exact expected result depends on the implemented business rule for referential integrity, which should be consistent).
    \end{itemize} \\ \hline

% Scenario 3: State Transition
\multicolumn{4}{|l|}{\textbf{Scenario: State Transitions Across Modules}} \\ \hline
TC\_INT\_004 & Verify that the state of a Book (isAvailable) is correctly updated by actions in the Borrowal module. & 
    \begin{itemize}[noitemsep, topsep=0pt, partopsep=0pt, leftmargin=*]
        \item \textbf{Prerequisite:} A book "StateTest Book" exists and is available. A "StateTest Member" exists.
        \item \textbf{As Member:} Login as "StateTest Member". Borrow "StateTest Book".
        \item \textbf{As Librarian:} Login. Navigate to Book Management. View the status of "StateTest Book".
        \item Try to create a new borrowal for "StateTest Book" for another member.
        \item Navigate to Borrowal Management. Find the borrowal record for "StateTest Book" and update its status to "Returned".
        \item Navigate back to Book Management and view the status of "StateTest Book" again.
    \end{itemize} & 
    \begin{itemize}[noitemsep, topsep=0pt, partopsep=0pt, leftmargin=*]
        \item After the member borrows the book, the book's `isAvailable` status immediately changes to `false`.
        \item The Librarian cannot create a new borrowal for the now-unavailable book. The UI should prevent this action.
        \item After the Librarian marks the borrowal as "Returned", the borrowal record's status field is updated.
        \item The `isAvailable` status of "StateTest Book" in the Book module reverts to `true`.
        \item The data remains consistent between the Book and Borrowal modules throughout the workflow.
    \end{itemize} \\ \hline

% Scenario 4: Role-Based Access Control
\multicolumn{4}{|l|}{\textbf{Scenario: Role-Based Access Control (RBAC) Integration}} \\ \hline
TC\_INT\_005 & Verify that user roles created in the User module correctly restrict access to features in other modules (e.g., Book, Author). &
    \begin{itemize}[noitemsep, topsep=0pt, partopsep=0pt, leftmargin=*]
        \item \textbf{As Librarian:} Login. Create a new user "TestMember" with the role of Member (`isAdmin=false`).
        \item Logout.
        \item \textbf{As Member:} Login as "TestMember".
        \item Attempt to access the User Management page via direct URL navigation (e.g., `/users`).
        \item Navigate to the Books list.
        \item Check for the presence of "Add New Book", "Edit", or "Delete" buttons/options.
        \item Using an API client (e.g., Postman), attempt to send a `POST` request to `/api/book/add` using the session/token of "TestMember".
    \end{itemize} & 
    \begin{itemize}[noitemsep, topsep=0pt, partopsep=0pt, leftmargin=*]
        \item The Member login is successful.
        \item Direct navigation to the `/users` URL is blocked, and the user is redirected to the member dashboard or an error page.
        \item The UI for the Books list does not show any administrative controls (Add, Edit, Delete).
        \item The API request to `POST /api/book/add` fails with an authorization error (e.g., 403 Forbidden or 401 Unauthorized), demonstrating that the access control middleware correctly uses the user's role from the session.
    \end{itemize} \\ \hline

\end{longtable}